\documentclass[12pt,a4paper]{article}

% Scientific Data template requirements
\usepackage{natbib}
\usepackage{graphicx}
\usepackage{hyperref}
\usepackage{url}
\usepackage{amsmath}
\usepackage{booktabs}
\usepackage{longtable}
\usepackage{array}
\usepackage{multirow}
\usepackage{wrapfig}
\usepackage{float}
\usepackage{colortbl}
\usepackage{pdflscape}
\usepackage{tabu}
\usepackage{threeparttable}
\usepackage{threeparttablex}
\usepackage{makecell}
\usepackage{xcolor}

% Metadata
\title{Darwin Operator Symmetry Atlas (DOSA): A Database of Dihedral Symmetries in Complete Bacterial Replicons}

\author{Demetrios Chiuratto Agourakis}

\date{\today}

\begin{document}

\maketitle

\begin{abstract}
We present the Darwin Operator Symmetry Atlas (DOSA), a comprehensive database of operator-defined symmetries in complete bacterial replicons. DOSA implements a hybrid Julia/Demetrios architecture to compute exact and approximate dihedral symmetry metrics across bacterial genomes. The atlas provides: (1) exact symmetry metrics including orbit sizes and palindrome detection; (2) approximate symmetry via normalized minimum dihedral distance (d\_min/L); (3) algebraic verification of dicyclic group structure (Dic\_n → D\_n double cover); and (4) an epistemic knowledge layer with full provenance tracking. All code and data are fully reproducible, with cross-validation ensuring identical results between implementations. The dataset is versioned and available via Zenodo DOI.
\end{abstract}

\section{Background \& Summary}

Bacterial genomes exhibit complex structural symmetries that reflect evolutionary constraints and functional requirements. The dihedral group D\_4, generated by cyclic shift (S) and reverse (R) operators, provides a mathematical framework for analyzing these symmetries. DOSA systematically computes and catalogs these symmetries across complete bacterial genomes from NCBI RefSeq.

The atlas addresses several key questions in bacterial genomics: (1) How prevalent are palindromic and reverse-complement fixed sequences? (2) What is the distribution of approximate symmetries across bacterial taxa? (3) How do symmetry patterns correlate with replication origin/terminus structure? (4) What is the enrichment of inverted repeats compared to random expectations?

Our dataset provides comprehensive answers to these questions through systematic computation of exact and approximate symmetry metrics, biological annotations (GC skew, ori/ter estimates, inverted repeats), and algebraic verification of group-theoretic structure.

\section{Methods}

\subsection{Data Acquisition}
Genomes were downloaded from NCBI RefSeq \cite{ncbi_refseq2023}, filtered for complete bacterial assemblies. The pipeline processes all replicons (chromosomes and plasmids) from each assembly. For the current release (v2.0.0-alpha), we analyzed 1,120 replicons from 500 complete bacterial genomes, sampled with seed=42 for reproducibility.

Each replicon is assigned a stable internal identifier (\texttt{replicon\_id}) and associated with metadata including assembly accession, replicon type, length, GC fraction, and taxonomy ID. SHA256 checksums ensure data integrity throughout the pipeline.

\subsection{Symmetry Computation}
\subsubsection{Exact Symmetry}
Orbit sizes under the dihedral group D\_4 are computed for each replicon. The dihedral group D\_4 consists of 4 elements: identity (I), reverse (R), complement (K), and reverse-complement (RC). For a sequence of length n, the orbit under D\_4 can have size 1, 2, or 4, corresponding to orbit ratios of 0.25, 0.5, or 1.0 respectively.

Fixed points under R (palindromes) and RC (reverse-complement fixed sequences) are identified. These represent sequences with maximal symmetry under the dihedral group.

\subsubsection{Approximate Symmetry}
Minimum dihedral distance d\_min is computed via exhaustive search over all 2n transforms (n cyclic shifts and n reverse-shifts), normalized by sequence length L to yield d\_min/L ∈ [0,1]. This metric quantifies how close a sequence is to exact dihedral symmetry, with d\_min/L = 0 indicating perfect symmetry.

\subsubsection{Algebraic Verification}
Dicyclic groups Dic\_n are verified to provide double covers of dihedral groups D\_n, with quaternion structure confirmed.

\subsection{Epistemic Knowledge Layer}
All metrics are exported with epistemic metadata: provenance (git SHA, timestamp, NCBI accession), uncertainty bounds (epsilon), confidence scores, and validity predicates.

\subsection{Reproducibility}
The pipeline is fully reproducible via \texttt{git clone} + \texttt{make reproduce}. Cross-validation ensures identical results between Julia (reference) and Demetrios (high-performance) implementations.

\section{Data Records}

The dataset (v2.0.0-alpha) includes 1,120 replicons from 500 complete bacterial genomes, with 80,659 epistemic knowledge records. All data is available in both CSV (for compatibility) and Parquet (for efficient querying) formats.

\subsection{Core Tables}

\begin{itemize}
\item \texttt{atlas\_replicons.csv} (1,120 rows): Replicon metadata including assembly accession (GCF format), length (bp), GC fraction, taxonomy ID, and SHA256 checksum.

\item \texttt{kmer\_inversion.csv} (11,190 rows): K-mer inversion symmetry metrics for k=1..10. Metrics include X\_k (asymmetry measure), K\_L(tau) (cumulative distribution), and replichore-specific values.

\item \texttt{gc\_skew\_ori\_ter.csv} (1,120 rows): GC skew analysis with ori/ter position estimates, confidence scores, and GC skew amplitude. Based on windowed GC skew computation similar to SkewDB \cite{skewdb2020}.

\item \texttt{replichore\_metrics.csv} (2,238 rows): Replichore-level metrics including length, GC fraction, and k-mer inversion (k=6) for leading and lagging strands.

\item \texttt{inverted\_repeats\_summary.csv} (1,120 rows): Inverted repeats detection with enrichment analysis. Includes IR count, density, baseline comparison, enrichment ratio, and statistical significance (p-value).

\item \texttt{quaternion\_results.csv} (15 rows): Dicyclic group verification results for n=2..16, confirming Dic\_n → D\_n double cover structure.

\item \texttt{dicyclic\_lifts.csv} (4 rows): Verification of dicyclic lifts for dihedral orders 4, 8, and 16.
\end{itemize}

\subsection{Epistemic Knowledge Layer}

The \texttt{epistemic/atlas\_knowledge.jsonl} file contains 80,659 knowledge records, each with:
\begin{itemize}
\item Provenance: Git SHA, timestamp, NCBI accession, replicon ID
\item Value: The computed metric
\item Epsilon: Error bound (0 for exact, null if unknown)
\item Confidence: Epistemic confidence in [0,1]
\item Validity: Domain constraint verification
\end{itemize}

This layer enables full traceability and uncertainty propagation in downstream analyses.

\section{Technical Validation}

\subsection{Cross-Validation}
Julia and Demetrios implementations produce identical results (535/535 tests passed).

\subsection{Data Integrity}
All checks passed:
\begin{itemize}
\item Schema compliance: 17/17 checks passed
\item Value ranges: All metrics within expected bounds
\item Referential integrity: All foreign keys valid
\item Epistemic validation: 1,124,386 checks passed, 0 failed
\end{itemize}

\subsection{Performance}
Demetrios FFI kernels show 1.5-9× speedup over Julia reference implementation, depending on sequence length. For \texttt{orbit\_size}, speedup ranges from 7.2× (short sequences, 16 bp) to 1.8× (very long sequences, 20,000 bp). For \texttt{dmin}, speedup ranges from 9.0× (short) to 1.5× (very long). The hybrid architecture provides both reproducibility (Julia reference) and performance (Demetrios kernels).

\section{Usage Notes}

\subsection{Reproducing the Dataset}
\begin{verbatim}
git clone https://github.com/chiuratto-AI/darwin-atlas.git
cd darwin-atlas
make atlas MAX=500 SEED=42
\end{verbatim}

\subsection{Querying the Dataset}
The dataset is available in both CSV and Parquet formats. SQL queries can be executed via DuckDB or similar tools. Parquet partitions are organized by source database and replicon type for efficient querying.

Example query to find replicons with high GC skew amplitude:
\begin{verbatim}
SELECT replicon_id, gc_skew_amplitude, ori_position, ter_position
FROM gc_skew_ori_ter
WHERE gc_skew_amplitude > 0.1
ORDER BY gc_skew_amplitude DESC
LIMIT 10;
\end{verbatim}

\section{Code Availability}

All code is available at: \url{https://github.com/chiuratto-AI/darwin-atlas}

\section{Data Availability}

The dataset is available via Zenodo DOI: [TO BE ASSIGNED]

\bibliographystyle{nature}
\bibliography{references}

\end{document}

